\section{Bugs from Specification Ambiguity}
\label{sec:bugs}

During KAT verification testing against the C reference implementation,
we discovered multiple bugs in our theta-coordinate $(2,2)$-isogeny
chain.  Each bug originated from an ambiguity in the specification or
an implicit convention in the reference code.  We catalog them here
with their root causes.

\subsection{ProductToTheta coordinate convention}
\label{sec:bug-product-to-theta}

The spec describes the product theta structure abstractly via
level-2 theta functions.  The relationship between Montgomery
projective coordinates $(X:Z)$ and the product theta representation
is not stated explicitly.  Our initial implementation used
$(X-Z, X+Z)$ as ``theta-like'' coordinates---natural for Montgomery
differential addition---but the C reference's \texttt{base\_change()}
uses raw $(X, Z)$ directly:
\begin{lstlisting}[language=C]
fp2_mul(&null.x, &T->P1.x, &T->P2.x); // X1 * X2
fp2_mul(&null.y, &T->P1.x, &T->P2.z); // X1 * Z2
fp2_mul(&null.z, &T->P2.x, &T->P1.z); // Z1 * X2
fp2_mul(&null.t, &T->P1.z, &T->P2.z); // Z1 * Z2
\end{lstlisting}

\noindent
\textbf{Lesson:} When the spec is abstract about coordinate
representations, the C reference is the ground truth.

\subsection{Squared theta means square then Hadamard}
\label{sec:bug-squared-theta}

In the SQIsign context, ``squared theta'' means component-wise
squaring \emph{followed by a Hadamard transform}.  The C reference
wraps this in \texttt{to\_squared\_theta()}, whose name does not
reveal the Hadamard.  Our generic isogeny evaluation (Algorithm~8.34)
computed $\mathcal{H}(\mathrm{inv} \cdot P^2)$ instead of
$\mathcal{H}(\mathrm{inv} \cdot \mathcal{H}(P^2))$.

\noindent
\textbf{Lesson:} Read the body of helper functions in the C reference,
not just their names.  ``Squared theta'' is not just squaring.

\subsection{Cross-product index errors in GluingEval}
\label{sec:bug-gluing-eval}

Algorithm~8.39 computes \texttt{ADDComponents} on each curve
component, producing $(u_1, v_1, w_1)$ from curve~1 and
$(u_2, v_2, w_2)$ from curve~2.  The $U$ vector's second entry
should be $u_1 \cdot w_2$ (cross-curve), but we wrote $u_1 \cdot w_1$
(same-curve).

\noindent
\textbf{Lesson:} In product-curve formulas, every multiplication
involves one factor from each curve.  Two factors with the same
subscript are almost certainly wrong.

\subsection{Projective inverse vs.\ field inversion}
\label{sec:bug-projective-inverse}

The gluing codomain needs the ``inverse'' of the dual theta null point
$(\alpha, \beta, \gamma, 0)$.  In projective coordinates this is
computed via cross-products---pure multiplications, no field inversion.
Our code called \texttt{Fp2::invert()}, which is both wrong and
expensive.

\noindent
\textbf{Lesson:} In projective or theta coordinates, ``inverse''
almost always means projective inverse via cross-products.  Actual
field inversion is a red flag.

\subsection{Reusing generic types for special patterns}
\label{sec:bug-imageK}

The image of the kernel generator under gluing has the pattern
$(x : x : y : y)$.  Storing this in a generic 4-coordinate
\texttt{JacobianPoint} creates ambiguity: \texttt{J.y} equals $x$
(the second copy), not $y$.  The scaling step used \texttt{J.x}
and \texttt{J.y} for the two distinct values, but both hold $x$.

\noindent
\textbf{Lesson:} Do not reuse a generic $n$-coordinate type to store
a value with a special pattern.  Use a dedicated type, or at minimum
document which fields hold which logical values.

\subsection{Torsion basis swap convention}
\label{sec:from-hint}

The C reference's \texttt{ec\_curve\_to\_basis\_2f\_from\_hint()}
deliberately stores $P-Q$ in \texttt{B.Q} and $Q$ in
\texttt{B.PmQ}:
\begin{lstlisting}[language=C]
difference_point(&PQ2->Q, &P, &Q, curve);  // B.Q = P - Q
copy_point(&PQ2->P, &P);                    // B.P = P
copy_point(&PQ2->PmQ, &Q);                  // B.PmQ = Q
\end{lstlisting}
This means the three-point ladder computes $P + [m](P-Q)$, not
$P + [m]Q$.  The swap ensures the multiplied point avoids $(0,0)$
for differential addition.  This is not documented in the spec.

\noindent
\textbf{Lesson:} Undocumented conventions in reference code can look
like bugs in your own code.  Verify against the reference before
``fixing'' something that isn't broken.


\subsection{Different formulas for the last two chain steps}
\label{sec:bug-hadamard-bool}

The C reference uses three different formula variants across the
$(2,2)$-isogeny chain, controlled by two boolean flags
\texttt{hadamard\_bool\_1} and \texttt{hadamard\_bool\_2}:
\begin{itemize}
  \item \textbf{Normal steps} ($i < n-2$): \texttt{bool\_1=0, bool\_2=1}.
    Codomain is Hadamard-transformed (standard form).
    Eval computes $\mathcal{H}(\mathrm{precomp} \cdot \mathcal{H}(P^2))$.
  \item \textbf{Penultimate step} ($i = n-2$): \texttt{bool\_1=0, bool\_2=0}.
    Codomain stays in dual form (no Hadamard).
    Eval computes $\mathrm{precomp} \cdot \mathcal{H}(P^2)$.
  \item \textbf{Ultimate step} ($i = n-1$): \texttt{bool\_1=1, bool\_2=0}.
    Input gets an extra Hadamard before squaring; codomain in dual form.
    Eval computes $\mathrm{precomp} \cdot \mathcal{H}(\mathcal{H}(P)^2)$.
\end{itemize}
The splitting step expects its input in dual form---what
\texttt{bool\_2=0} produces.  Our code used the normal-step formula
(\texttt{bool\_2=1}) for all steps, feeding the splitting a
Hadamard-transformed null point that has no zero $U_{i,j}(0)$ entry.
The chain ran correctly for 125 steps and produced nonsense at the end.

\noindent
\textbf{Lesson:} The last few steps of an isogeny chain often need
special handling.  The spec describes the generic step uniformly; the
\texttt{hadamard\_bool} optimization is an implementation-level choice
in the C reference that avoids redundant Hadamard transforms at the
boundary between the chain and the splitting step.  Check the loop
structure, not just the per-step formula.

\subsection{\texorpdfstring{$\mathbb{F}_{p^2}$}{Fp2} square root canonicalization}
\label{sec:bug-sqrt}

This was the single most impactful bug we encountered, and we devote
extra space to it because its consequences are subtle and far-reaching.

Every nonzero square in $\mathbb{F}_{p^2}$ has exactly two roots,
$r$ and $-r$.  The SQIsign specification (Algorithm~8.12, the
\texttt{difference\_point} subroutine of \texttt{TorsionBasisFromHint})
says ``compute a square root'' without specifying which one.  We
initially returned whichever root the Tonelli--Shanks formula
produced.

The C reference's \texttt{fp2\_sqrt} (described in Aardal et
al.~\cite{ABCDK24}, \S3.1) has a five-line canonicalization step at
the end: it negates the output whenever the real part is odd (as an
integer in $[0, p)$), or when the real part is zero and the imaginary
part is odd.  This ensures a unique, deterministic ``even'' root for
every input.

Without this canonicalization, \texttt{projective\_difference}
nondeterministically adds $+\sqrt{\Delta}$ or $-\sqrt{\Delta}$ to
$B_{XZ}$, producing two affinely distinct candidates for
$x(P-Q)$.  Only one matches the C reference's
\texttt{difference\_point}.  The wrong choice propagates through
\texttt{ladder3pt} (producing a different kernel generator), the
challenge isogeny (landing on a different curve $E_{\text{chl}}$
with a different $j$-invariant), and every subsequent step.  In our
case, the challenge curve's $j$-invariant was $\mathtt{0x03007c\ldots}$
instead of the correct $\mathtt{0x026558\ldots}$---completely unrelated
values, not just a sign flip.

The key insight is that the choice of square root is \emph{not absorbed
by projective equivalence}.  The formula
$x_{P-Q} = (\sqrt{\Delta} + B_{XZ} \;:\; B_{ZZ})$
is not projectively equivalent to
$(-\sqrt{\Delta} + B_{XZ} \;:\; B_{ZZ})$;
these are genuinely different points.  Any implementation of SQIsign
that computes torsion bases via \texttt{difference\_point} must
canonicalize its $\mathbb{F}_{p^2}$ square root.

\noindent
\textbf{Lesson:}  Any time a specification says ``compute a square
root'' in $\mathbb{F}_{p^2}$, the implementation \emph{must} specify
and enforce a canonical choice.  The even-real-part convention used by
the C reference is simple and should be the default for any SQIsign
implementation.  This is not documented in the SQIsign specification
as of v2.0.1.

\subsection{Never recompute \texorpdfstring{$P{-}Q$}{P-Q} via
  \texttt{DifferencePoint}}
\label{sec:bug-pmq-recompute}

A torsion basis $(P, Q, P{-}Q)$ is produced once by
\texttt{TorsionBasisFromHint}.  The third component $P{-}Q$ is
computed via \texttt{DifferencePoint}, which takes an
$\mathbb{F}_{p^2}$ square root.  We discovered---after fixing the
sqrt canonicalization---that $P{-}Q$ must \emph{never} be recomputed
from $P$ and $Q$ in any subsequent operation.  Even with a canonical
sqrt, the two roots of \texttt{DifferencePoint} yield the affinely
distinct points $x(P{-}Q)$ and $x(2P{-}Q)$, and there is no guarantee
that a fresh call returns the same one.

We found three separate sites where recomputation caused failure:
\begin{enumerate}
  \item \textbf{After scaling.} We doubled $P$ and $Q$ to reduce
    their order, then called \texttt{projective\_difference} on the
    doubled points instead of doubling the original $P{-}Q$ alongside.
  \item \textbf{In the matrix application.} $M_{\mathrm{chl}} \cdot
    (P, Q)$ produced $P'$ and $Q'$ via two biladder calls.  We
    computed $P'{-}Q'$ via \texttt{projective\_difference} instead of
    a third biladder with scalars $(a{-}b, c{-}d)$, as the C reference
    does in \texttt{matrix\_scalar\_application\_even\_basis}.
  \item \textbf{Before the $(2,2)$-chain.} After the even response
    isogeny, we recomputed $P{-}Q$ from the images instead of pushing
    the original $P{-}Q$ through the isogeny as a third evaluation
    point.
\end{enumerate}
Each fix was individually necessary; all three together were sufficient
to make KAT vector~0 verify.

\noindent
\textbf{Lesson:} In $x$-only Montgomery arithmetic, $P{-}Q$ is a
\emph{derived value} that cannot be reconstructed from $x(P)$ and
$x(Q)$ alone (because $x(P{-}Q) = x(P{+}Q)$ on the Kummer line, and
\texttt{DifferencePoint} resolves this ambiguity via a square root
whose sign is fragile).  Treat $P{-}Q$ as a first-class member of the
basis triple and propagate it through every transformation.

